\hypertarget{customer}{%
\section{Customer}\label{customer}}

\hypertarget{configure-bash-for-xml-api}{%
\subsection{Configure BASH for XML
API}\label{configure-bash-for-xml-api}}

Export these environment variables to prepare for script execution.

NOTE: You may also specify USER and PAN\_IP by passing flags to the
script.

\begin{Shaded}
\begin{Highlighting}[]
\BuiltInTok{export} \VariableTok{USER=}\OperatorTok{\textless{}}\ExtensionTok{the\_name\_of\_the\_firewall\_administrator\_account\_you\_configured}\OperatorTok{\textgreater{}}\NormalTok{ \# OPTIONAL}
\BuiltInTok{export} \VariableTok{PASS=}\OperatorTok{\textless{}}\ExtensionTok{the\_password\_of\_the\_firewall\_administrator\_account\_you\_configured}\OperatorTok{\textgreater{}}
\BuiltInTok{export} \VariableTok{PAN\_IP=}\OperatorTok{\textless{}}\ExtensionTok{the\_panorama\_instance\_ip\_address}\OperatorTok{\textgreater{}}\NormalTok{ \# OPTIONAL}
\end{Highlighting}
\end{Shaded}

\hypertarget{optional-bash-env-var-user}{%
\subsubsection{OPTIONAL BASH ENV VAR:
USER}\label{optional-bash-env-var-user}}

This is the username you created previously for the dedicated XML API
user.

NOTE: You may also specify a user from the command line with the
\texttt{-u} flag.

For the \texttt{USER} setting, the tool has a default value of
\texttt{xml-api-user}. To override this value, use the following command
before running the tool:

\begin{Shaded}
\begin{Highlighting}[]
\BuiltInTok{export} \VariableTok{USER=}\StringTok{"my\_new\_username"}\NormalTok{ \# }\VariableTok{set} \VariableTok{the} \VariableTok{value}
\BuiltInTok{echo} \VariableTok{$USER} \CommentTok{\# validate it was set properly}
\end{Highlighting}
\end{Shaded}

\hypertarget{bash-environment-variable-pass}{%
\subsubsection{BASH ENVIRONMENT VARIABLE:
PASS}\label{bash-environment-variable-pass}}

This is the password for the new XML API user that was created. It is
important that we do not hardcode any active passwords into the script.

For the \texttt{PASS} setting, the tool has a default value of
\texttt{default}. To override this value, use the following command
before running the tool:

\begin{Shaded}
\begin{Highlighting}[]
\BuiltInTok{export} \VariableTok{PASS=}\StringTok{"my\_secret password"}\NormalTok{ \# }\VariableTok{set} \VariableTok{the} \VariableTok{value}
\BuiltInTok{echo} \VariableTok{$PASS} \CommentTok{\# validate it was set properly}
\KeywordTok{set} \KeywordTok{|} \FunctionTok{grep} \VariableTok{$PASS} \CommentTok{\# different way to validate}
\end{Highlighting}
\end{Shaded}

Alternatively, you may choose to use a combination of the
\href{https://direnv.net/}{direnv} and
\href{https://www.passwordstore.org/}{pass} tools to make the management
of these variables slightly easier.

\hypertarget{optional-bash-env-var-pan_ip}{%
\subsubsection{OPTIONAL BASH ENV VAR:
PAN\_IP}\label{optional-bash-env-var-pan_ip}}

This is the IPv4 address of the active Panorama.

NOTE:

\begin{itemize}
\tightlist
\item
  You may specify TESTNET Panorama IP from the command line with the
  \texttt{-t} flag.
\item
  You may specify PRODUCTION Panorama IP from the command line with the
  \texttt{-p} flag.
\end{itemize}

For the \texttt{PAN\_IP} setting, the tool has a default value of
\texttt{34.134.31.136}. To override this value, use the following
command before running the tool:

\begin{Shaded}
\begin{Highlighting}[]
\BuiltInTok{export} \VariableTok{PAN\_IP=}\StringTok{"10.2.3.4"}\NormalTok{ \# }\VariableTok{example}\NormalTok{, }\VariableTok{set} \VariableTok{the} \VariableTok{value}
\BuiltInTok{echo} \VariableTok{$PAN\_IP} \CommentTok{\# validate it was set properly}
\KeywordTok{set} \KeywordTok{|} \FunctionTok{grep} \VariableTok{$PAN\_IP} \CommentTok{\# another way to validate}
\end{Highlighting}
\end{Shaded}

\hypertarget{panorama-xml-api}{%
\subsection{Panorama XML API}\label{panorama-xml-api}}

In order to allow the BASH script to communicate with Panorama via the
XML API, you need to
\href{https://docs.paloaltonetworks.com/pan-os/10-2/pan-os-panorama-api/get-started-with-the-pan-os-xml-api/enable-api-access.html}{enable
the API access}.

\textbf{As best practice, set up a separate admin account for the XML
API access.}

\begin{enumerate}
\def\labelenumi{\arabic{enumi}.}
\item
  Configure an Admin Role Profile:

  \begin{itemize}
  \tightlist
  \item
    Select \texttt{Device} \textgreater{} \texttt{Admin\ Roles} and
    click \texttt{Add}.
  \item
    Enter a \texttt{Name} to identify the role.
  \item
    Select the \texttt{Role} scope.
  \item
    In the \texttt{XML\ API} tab, click the icon for each functional
    area to toggle it to the desired setting: \texttt{Enable} or
    \texttt{Disable}.
  \item
    Click \texttt{OK} to save the profile.
  \item
    Assign the role to an administrator - complete the
    \texttt{Configure\ a\ Firewall\ Administrator\ Account} step.
  \end{itemize}
\item
  Configure a Firewall Administrator Account:

  \begin{itemize}
  \tightlist
  \item
    Select \texttt{Device} \textgreater{} \texttt{Administrators} and
    click \texttt{Add}.
  \item
    Enter a user \texttt{Name} to identify the account. If the firewall
    uses a local user database to authenticate the account, enter the
    name that you specified for the account in the database.
  \item
    Select an \texttt{Authentication\ Profile} or sequence if you
    configured either for the administrator. If the firewall uses
    \texttt{Local\ Authentication} without a local user database for the
    account, select \texttt{None} (default) and enter a
    \texttt{Password}.
  \item
    Select the \texttt{Administrator\ Type} \textgreater{}
    \texttt{Custom\ Panorama\ Admin}.
  \item
    Select the \texttt{Admin\ Role} you created for this account in the
    \texttt{Profile} section.
  \item
    Click \texttt{OK} and \texttt{Commit}.
  \end{itemize}
\end{enumerate}

NOTE: You must create the admin role and XML user in two places. The
Panorama and the FW template/template stack must both have the role and
dedicated XML user added.

\hypertarget{validate-xml-api-configuration}{%
\subsubsection{Validate XML API
configuration}\label{validate-xml-api-configuration}}

To check if the \texttt{XML\ API} has been enabled, run the following
code snippet:

\begin{Shaded}
\begin{Highlighting}[]
\ExtensionTok{curl}\NormalTok{ {-}k {-}X GET }\StringTok{\textquotesingle{}https://$PAN\_IP//api/?type=keygen\&user=$USER\&password=$PASS\textquotesingle{}}
\end{Highlighting}
\end{Shaded}

If everything was configured correctly, a response status
\texttt{success} should be returned together with an API key for the
specified account.

\hypertarget{cli-dev}{%
\subsubsection{CLI Dev}\label{cli-dev}}

\begin{itemize}
\tightlist
\item
  \href{https://docs.paloaltonetworks.com/pan-os/11-0/pan-os-panorama-api/get-started-with-the-pan-os-xml-api/explore-the-api/use-the-cli-to-find-xml-api-syntax}{Use
  the CLI to Find XML API Syntax}
\end{itemize}

\hypertarget{script-binstatus.sh}{%
\subsection{\texorpdfstring{SCRIPT:
\texttt{bin/status.sh}}{SCRIPT: bin/status.sh}}\label{script-binstatus.sh}}

\hypertarget{disk-status}{%
\subsubsection{Disk Status}\label{disk-status}}

The script should be updated to check disk status of log collectors.

\begin{Shaded}
\begin{Highlighting}[]
\FunctionTok{tail}\NormalTok{ follow yes  mp{-}log vld{-}1{-}0.log}
\FunctionTok{tail}\NormalTok{ follow yes  mp{-}log vld{-}2{-}0.log}
\ExtensionTok{show}\NormalTok{ system disk{-}space}
\FunctionTok{less}\NormalTok{ mp{-}log vldmgr.log}
\FunctionTok{less}\NormalTok{ mp{-}log logd.log}
\FunctionTok{less}\NormalTok{ mp{-}log configd.log }\KeywordTok{|} \ExtensionTok{match}\NormalTok{ Error}
\ExtensionTok{show}\NormalTok{ logging{-}status all}
\ExtensionTok{show}\NormalTok{ system software status }\KeywordTok{|} \ExtensionTok{match} \StringTok{"elasticsearch\textbackslash{}|es{-}"}
\ExtensionTok{show}\NormalTok{ system raid detail}
\ExtensionTok{debug}\NormalTok{ log{-}collector log{-}collection{-}stats show incoming{-}logs}
\ExtensionTok{request}\NormalTok{ log{-}fwd{-}ctrl device 010108010490 action start{-}from{-}lastack}
\end{Highlighting}
\end{Shaded}

\hypertarget{script-binstage_release}{%
\subsection{\texorpdfstring{SCRIPT:
\texttt{bin/stage\_release}}{SCRIPT: bin/stage\_release}}\label{script-binstage_release}}

The following list can be generated by calling the script with the
\texttt{-h} flag.

\begin{Shaded}
\begin{Highlighting}[]
\ExtensionTok{franklin@ups}\NormalTok{:\textasciitilde{}/bin$ ./stage\_release.sh {-}h}

\ExtensionTok{stage\_release}\NormalTok{ script.}

\ExtensionTok{Syntax}\NormalTok{: stage\_release.sh [{-}f}\KeywordTok{|}\ExtensionTok{{-}h}\KeywordTok{|}\ExtensionTok{{-}i}\KeywordTok{|}\ExtensionTok{{-}p}\KeywordTok{|}\ExtensionTok{{-}r}\KeywordTok{|}\ExtensionTok{{-}t}\KeywordTok{|}\ExtensionTok{{-}u}\KeywordTok{|}\ExtensionTok{{-}v}\NormalTok{]}
\ExtensionTok{options}\NormalTok{:}
\ExtensionTok{{-}f}\NormalTok{     Specify a file with FW IPv4 addresses, one per line.}
\ExtensionTok{{-}h}\NormalTok{     Print this Help.}
\ExtensionTok{{-}i}\NormalTok{     Install release}
\ExtensionTok{{-}p}\NormalTok{     Use PROD environment}
\ExtensionTok{{-}r}\NormalTok{     Specify a sofwtware release to download/install (ex: 10.2.0)}
\ExtensionTok{{-}t}\NormalTok{     Use TESTNET environment}
\ExtensionTok{{-}u}\NormalTok{     Specify a USER for the Panorama}
\ExtensionTok{{-}v}\NormalTok{     Enable verbose (debug) }\ExtensionTok{mode.}
\end{Highlighting}
\end{Shaded}

\hypertarget{execution}{%
\subsubsection{Execution}\label{execution}}

Create a text file and add in IP addresses, one per line. This will
allow you to specify exactly which firewalls to operate on.

Example:

\begin{Shaded}
\begin{Highlighting}[]
\ExtensionTok{3.233.53.199}
\ExtensionTok{44.216.25.244}
\ExtensionTok{34.206.152.182}
\ExtensionTok{52.55.185.160}
\end{Highlighting}
\end{Shaded}

To run the script with a list of serials in a file, type the command
\texttt{./stage\_release.sh\ -f\ my\_ip\_list.txt}

To run the script against all firewalls that are managed by this
Panorama, type the command \texttt{stage\_release.sh}

\hypertarget{dev-notes}{%
\subsubsection{Dev Notes}\label{dev-notes}}

\begin{itemize}
\tightlist
\item
  Use \texttt{pandoc} to generate documentation
\end{itemize}

\begin{Shaded}
\begin{Highlighting}[]
\FunctionTok{sudo}\NormalTok{ apt install pandoc}
\BuiltInTok{cd}\NormalTok{ docs }\KeywordTok{\&\&} \ExtensionTok{pandoc}\NormalTok{ stage\_release.md {-}o stage\_release.pdf}
\end{Highlighting}
\end{Shaded}

\hypertarget{jump-host-setup}{%
\subsection{Jump Host Setup}\label{jump-host-setup}}

RHEL/CentOS

\begin{Shaded}
\begin{Highlighting}[]
\FunctionTok{sudo}\NormalTok{ yum install {-}y epel{-}release pandoc git}
\FunctionTok{sudo}\NormalTok{ yum install {-}y gpg2 }\CommentTok{\# for signing commits, encrypting passwords, etc.}
\FunctionTok{sudo}\NormalTok{ yum install {-}y direnv }\CommentTok{\# manage env vars}
\end{Highlighting}
\end{Shaded}

\hypertarget{pass}{%
\subsubsection{pass}\label{pass}}

the \texttt{pass} command as described in this article:
\url{https://www.redhat.com/sysadmin/management-password-store}

\hypertarget{development-environment}{%
\subsubsection{development environment}\label{development-environment}}

Run the command: sudo yum group install ``Development Tools'' Because I
would like to configure the development env as detailed here: CentOS /
RHEL 7: Install GCC (C and C++ Compiler) and Development Tools -
nixCraft (cyberciti.biz)

\hypertarget{services}{%
\subsubsection{Services}\label{services}}

\begin{itemize}
\tightlist
\item
  nginx
\item
  postfix/sendmail
\item
  crontab
\end{itemize}

\hypertarget{google-cloud}{%
\subsection{Google Cloud}\label{google-cloud}}

\hypertarget{gcp-production}{%
\subsubsection{GCP Production}\label{gcp-production}}

\includegraphics{docs/drawings/diagram1-GCP-PROD.png}

\hypertarget{gcp-testnet}{%
\subsubsection{GCP Testnet}\label{gcp-testnet}}

\includegraphics{docs/drawings/diagram1-GCP-TESTNET.png}

\hypertarget{sso-testing-in-gcp}{%
\subsubsection{SSO Testing in GCP}\label{sso-testing-in-gcp}}

\includegraphics{docs/drawings/diagram1-sandbox.png}

\hypertarget{prod-and-testnet}{%
\subsection{PROD and TESTNET}\label{prod-and-testnet}}

\hypertarget{new-jersey}{%
\subsubsection{New Jersey}\label{new-jersey}}

\includegraphics{docs/drawings/diagram1-New Jersey.png}

\hypertarget{georgia}{%
\subsubsection{Georgia}\label{georgia}}

\includegraphics{docs/drawings/diagram1-Georgia.png}

\hypertarget{github}{%
\subsection{Github}\label{github}}

\hypertarget{franklin}{%
\subsubsection{Franklin}\label{franklin}}

\hypertarget{customer-internal}{%
\subsubsection{customer internal}\label{customer-internal}}
